\chapter{§2.2 实数}
\label{sec:2.2}


\prerequisite{§2.1 有理数}

\gradelevel{8 年级}


\section*{平方根与算术平方根}


\begin{explore}


学校要在操场上铺设一块正方形的草坪，面积恰好为 $25$ 平方米。正方形的边长是多少？很简单： $5 \times 5 = 25$ ，所以边长是 $5$ 米。那如果面积是 $9$ 平方米呢？边长是 $3$ 米。但如果面积是 $2$ 平方米，边长应该是多少？我们找不到一个有理数使得它的平方等于 $2$ ——这就引出了"平方根"的概念。

\end{explore}

\begin{understand}


\textbf{平方根}：如果 $x^2 = a$ （ $a \geq 0$ ），则 $x$ 叫做 $a$ 的\textbf{平方根}（也叫二次方根）。

\begin{itemize}
  \item $a > 0$ 时， $a$ 有两个平方根： $\sqrt{a}$ 和 $-\sqrt{a}$ ，记为 $\pm\sqrt{a}$ 。
  \item $a = 0$ 时， $0$ 的平方根是 $0$ 。
  \item $a < 0$ 时， $a$ 没有平方根（在实数范围内）。
\end{itemize}


\textbf{算术平方根}：非负数 $a$ 的\textbf{非负平方根}叫做 $a$ 的算术平方根，记为 $\sqrt{a}$ 。

\[
\sqrt{a} \geq 0 \quad (a \geq 0)
\]


\textbf{为什么要区分"平方根"和"算术平方根"？} $9$ 的平方根有两个： $3$ 和 $-3$ 。但当我们写 $\sqrt{9}$ 时，特指非负的那个，即 $\sqrt{9} = 3$ 。这种约定避免了歧义。

\end{understand}

\begin{workedexamples}


\textbf{例 1.} 求下列各数的平方根： $16$ ， $\frac{9}{4}$ ， $0$ 。

\textbf{解：}

\begin{itemize}
  \item $16$ 的平方根：因为 $(\pm 4)^2 = 16$ ，所以 $16$ 的平方根为 $\pm 4$ 。
  \item $\frac{9}{4}$ 的平方根：因为 $\left(\pm \frac{3}{2}\right)^2 = \frac{9}{4}$ ，所以 $\frac{9}{4}$ 的平方根为 $\pm \frac{3}{2}$ 。
  \item $0$ 的平方根： $0^2 = 0$ ，所以 $0$ 的平方根为 $0$ 。
\end{itemize}


\textbf{例 2.} 求 $\sqrt{49}$ ， $-\sqrt{0.81}$ ， $\sqrt{(-3)^2}$ 。

\textbf{解：}

\begin{itemize}
  \item $\sqrt{49} = 7$ （算术平方根取非负值）。
  \item $-\sqrt{0.81} = -0.9$ （先求算术平方根 $0.9$ ，再取负号）。
  \item $\sqrt{(-3)^2} = \sqrt{9} = 3$ （注意：结果不是 $-3$ ，因为算术平方根非负）。
\end{itemize}


\textbf{例 3.} 已知一个正方形的面积为 $50$ 平方厘米，求其边长。

\textbf{解：} 设边长为 $a$ 厘米，则 $a^2 = 50$ ， $a = \sqrt{50}$ 。

\[
\sqrt{50} = \sqrt{25 \times 2} = 5\sqrt{2} \approx 7.07 \text{（厘米）}
\]


\end{workedexamples}

\begin{keytakeaway}


\begin{itemize}
  \item 正数有两个平方根（一正一负），零的平方根是零，负数没有实数平方根。
  \item $\sqrt{a}$ 表示算术平方根，永远非负。
  \item $\sqrt{a^2} = \lvert a \rvert$ ，不是 $a$ 。
\end{itemize}


\end{keytakeaway}

\begin{practice}


\begin{enumerate}
  \item 求 $64$ 的平方根。
  \item 求 $\sqrt{121}$ 。
  \item $-\sqrt{36}$ 等于多少？
  \item $\sqrt{(-5)^2}$ 等于多少？
  \item 判断对错： $\sqrt{16} = \pm 4$ 。
  \item 面积为 $200$ 平方米的正方形，边长是多少（用根号表示）？
\end{enumerate}


\end{practice}

\section*{立方根}


\begin{explore}


一个正方体的体积为 $8$ 立方厘米，它的棱长是多少？因为 $2 \times 2 \times 2 = 8$ ，所以棱长是 $2$ 厘米。如果体积是 $27$ 立方厘米呢？棱长是 $3$ 厘米。如果体积是 $-8$ 立方厘米（理论上，负体积可以表示"挖去"的空间），对应的"棱长"该是多少？

\end{explore}

\begin{understand}


\textbf{立方根}（三次方根）：如果 $x^3 = a$ ，则 $x$ 叫做 $a$ 的\textbf{立方根}，记为 $\sqrt[3]{a}$ 。

与平方根不同的是：

\begin{itemize}
  \item 每个实数（正数、零、负数）\textbf{都有且只有一个}立方根。
  \item 正数的立方根为正，零的立方根为零，负数的立方根为负。
\end{itemize}


\[
\sqrt[3]{8} = 2, \quad \sqrt[3]{0} = 0, \quad \sqrt[3]{-27} = -3
\]


\textbf{为什么负数有立方根而没有（实数）平方根？} 因为一个负数的立方仍然是负数（ $(-3)^3 = -27$ ），但任何实数的平方都是非负的（ $(-3)^2 = 9 > 0$ ）。

\end{understand}

\begin{workedexamples}


\textbf{例 1.} 求下列各数的立方根： $125$ ， $-64$ ， $\frac{1}{8}$ 。

\textbf{解：}

\begin{itemize}
  \item $\sqrt[3]{125} = 5$ ，因为 $5^3 = 125$ 。
  \item $\sqrt[3]{-64} = -4$ ，因为 $(-4)^3 = -64$ 。
  \item $\sqrt[3]{\frac{1}{8}} = \frac{1}{2}$ ，因为 $\left(\frac{1}{2}\right)^3 = \frac{1}{8}$ 。
\end{itemize}


\textbf{例 2.} 求 $\sqrt[3]{-0.001}$ 。

\textbf{解：} $(-0.1)^3 = -0.001$ ，所以 $\sqrt[3]{-0.001} = -0.1$ 。

\textbf{例 3.} 一个正方体容器的容积为 $216$ 立方厘米，求棱长。

\textbf{解：} 设棱长为 $a$ ，则 $a^3 = 216$ ， $a = \sqrt[3]{216} = 6$ 。棱长为 $6$ 厘米。

\end{workedexamples}

\begin{keytakeaway}


\begin{itemize}
  \item 每个实数都恰好有一个立方根。
  \item 立方根的符号与原数相同。
  \item $\sqrt[3]{a^3} = a$ （对所有实数 $a$ 成立，无需加绝对值）。
\end{itemize}


\end{keytakeaway}

\begin{practice}


\begin{enumerate}
  \item 求 $\sqrt[3]{343}$ 。
  \item 求 $\sqrt[3]{-1000}$ 。
  \item 求 $\sqrt[3]{\frac{27}{8}}$ 。
  \item 求 $\sqrt[3]{-0.125}$ 。
  \item 判断对错：负数没有立方根。
  \item 一个正方体的体积为 $1728$ 立方厘米，求棱长。
\end{enumerate}


\end{practice}

\section*{无理数的发现}


\begin{explore}


古希腊数学家毕达哥拉斯相信"万物皆数"——他认为世界上所有的量都可以用整数之比（即有理数）来表示。然而，他的学生希帕索斯发现：边长为 $1$ 的正方形，其对角线的长度无法用任何分数表示！这个长度就是 $\sqrt{2}$ 。这个发现在当时引起了极大的震动。

\end{explore}

\begin{understand}


\textbf{无理数}：不能写成两个整数之比 $\frac{p}{q}$ （ $q \neq 0$ ）的数，叫做无理数。

常见的无理数： $\sqrt{2}$ ， $\sqrt{3}$ ， $\pi$ ， $\sqrt[3]{2}$ 。

无理数的小数形式是\textbf{无限不循环小数}。例如：

\[
\sqrt{2} = 1.41421356\ldots \quad (\text{不循环})
\]


相反，有理数的小数形式要么是有限小数，要么是无限循环小数。

\textbf{ $\sqrt{2}$ 为什么不是有理数？（选读）}

用反证法：假设 $\sqrt{2} = \frac{p}{q}$ ，其中 $p, q$ 为互质的正整数。两边平方得 $2 = \frac{p^2}{q^2}$ ，即 $p^2 = 2q^2$ 。这说明 $p^2$ 是偶数，所以 $p$ 是偶数。设 $p = 2k$ ，代入得 $4k^2 = 2q^2$ ，即 $q^2 = 2k^2$ ，于是 $q$ 也是偶数。但这与 $p, q$ 互质矛盾！所以 $\sqrt{2}$ 不是有理数。

\end{understand}

\begin{workedexamples}


\textbf{例 1.} 判断下列各数是有理数还是无理数：

(1) $\sqrt{9}$ 　　(2) $\sqrt{5}$ 　　(3) $3.14$ 　　(4) $\pi$

\textbf{解：}

(1) $\sqrt{9} = 3$ ，是有理数。

(2) $\sqrt{5} = 2.2360679\ldots$ ，不能写成分数，是无理数。

(3) $3.14$ 是有限小数，可以写成 $\frac{314}{100}$ ，是有理数。

(4) $\pi = 3.14159265\ldots$ ，是无限不循环小数，是无理数。注意 $3.14$ 只是 $\pi$ 的近似值。

\textbf{例 2.} 写出一个大于 $3$ 小于 $4$ 的无理数。

\textbf{解：} $\sqrt{10}$ ，因为 $3^2 = 9 < 10 < 16 = 4^2$ ，所以 $3 < \sqrt{10} < 4$ 。

\textbf{例 3.} 估计 $\sqrt{7}$ 的整数部分。

\textbf{解：} 因为 $2^2 = 4 < 7 < 9 = 3^2$ ，所以 $2 < \sqrt{7} < 3$ ， $\sqrt{7}$ 的整数部分为 $2$ 。

\end{workedexamples}

\begin{keytakeaway}


\begin{itemize}
  \item 无理数是无限不循环小数，不能写成分数。
  \item $\sqrt{n}$ （ $n$ 为正整数）当 $n$ 不是完全平方数时为无理数。
  \item $\pi$ 是无理数， $3.14$ 只是近似值。
\end{itemize}


\end{keytakeaway}

\begin{practice}


\begin{enumerate}
  \item 判断 $\sqrt{36}$ 是有理数还是无理数。
  \item 判断 $\sqrt{11}$ 是有理数还是无理数。
  \item 写出一个大于 $1$ 小于 $2$ 的无理数。
  \item 估计 $\sqrt{50}$ 的整数部分。
  \item 判断对错：无限小数一定是无理数。
  \item 下列数中哪些是无理数： $0.333\ldots$ ， $\sqrt{8}$ ， $-\frac{22}{7}$ ， $\pi - 3$ ？
\end{enumerate}


\end{practice}

\section*{实数的分类与实数轴}


\begin{explore}


有了有理数和无理数，我们的数的世界变得更加完整了。把所有有理数和无理数放在一起，就构成了\textbf{实数}。一个自然的问题是：数轴上是否每个点都对应一个实数？答案是肯定的——这就是实数轴的完备性。

\end{explore}

\begin{understand}


\textbf{实数的分类}：

\[
\text{实数}\begin{cases}
\text{有理数}\begin{cases} \text{整数} \\ \text{分数} \end{cases} \\[6pt]
\text{无理数}
\end{cases}
\]


也可以按正负分：

\[
\text{实数}\begin{cases}
\text{正实数} \\
0 \\
\text{负实数}
\end{cases}
\]


\textbf{实数轴}：实数与数轴上的点\textbf{一一对应}。

\begin{itemize}
  \item 有理数可以在数轴上精确标出。
  \item 无理数也对应数轴上确定的点。例如 $\sqrt{2}$ 对应边长为 $1$ 的正方形对角线长度，可以用圆规在数轴上画出。
\end{itemize}


\textbf{实数的运算性质}：有理数的运算律（交换律、结合律、分配律）对所有实数仍然成立。相反数、绝对值的定义也自然推广到实数。

\end{understand}

\begin{workedexamples}


\textbf{例 1.} 把下列各数分别归入有理数集合与无理数集合：

\[
-3,\; \sqrt{2},\; 0,\; \frac{22}{7},\; \pi,\; -\sqrt{16},\; 0.\dot{3}
\]


\textbf{解：}

\begin{itemize}
  \item 有理数集合： $\{-3,\; 0,\; \frac{22}{7},\; -\sqrt{16},\; 0.\dot{3}\}$
\end{itemize}


  （ $-\sqrt{16} = -4$ ， $0.\dot{3} = \frac{1}{3}$ ， $\frac{22}{7}$ 是分数）

\begin{itemize}
  \item 无理数集合： $\{\sqrt{2},\; \pi\}$
\end{itemize}


\textbf{例 2.} 在数轴上， $\sqrt{3}$ 大约在什么位置？

\textbf{解：} $1^2 = 1 < 3 < 4 = 2^2$ ，所以 $1 < \sqrt{3} < 2$ 。更精确地， $1.7^2 = 2.89$ ， $1.8^2 = 3.24$ ，所以 $\sqrt{3} \approx 1.73$ ，在数轴上 $1$ 和 $2$ 之间靠近 $1.7$ 的位置。

\textbf{例 3.} 比较 $\sqrt{5}$ 与 $2.3$ 的大小。

\textbf{解：} $\sqrt{5} \approx 2.236$ ， $2.3 = 2.300$ 。因为 $2.236 < 2.300$ ，所以 $\sqrt{5} < 2.3$ 。

也可以用平方比较： $(\sqrt{5})^2 = 5$ ， $2.3^2 = 5.29$ 。因为 $5 < 5.29$ ，所以 $\sqrt{5} < 2.3$ 。

\textbf{例 4.} 求 $\lvert 2 - \sqrt{5} \rvert$ 。

\textbf{解：} 因为 $\sqrt{5} \approx 2.236 > 2$ ，所以 $2 - \sqrt{5} < 0$ 。

\[
\lvert 2 - \sqrt{5} \rvert = -(2 - \sqrt{5}) = \sqrt{5} - 2
\]


\end{workedexamples}

\begin{keytakeaway}


\begin{itemize}
  \item 实数 = 有理数 + 无理数。
  \item 实数与数轴上的点一一对应（实数轴的完备性）。
  \item 比较含根号的实数大小，可以利用平方法。
  \item 化简含绝对值的式子时，先判断被绝对值号包裹的表达式的正负。
\end{itemize}


\end{keytakeaway}

\begin{practice}


\begin{enumerate}
  \item 将 $\sqrt{25}$ 、 $\sqrt{7}$ 、 $-\pi$ 、 $\frac{1}{3}$ 、 $0$ 分为有理数和无理数。
  \item 比较大小： $\sqrt{10}$ 与 $3.2$ 。
  \item 求 $\lvert \sqrt{3} - 2 \rvert$ 。
  \item 估计 $\sqrt{20}$ 的值（精确到十分位）。
  \item 判断对错：所有实数都能在数轴上找到对应的点。
  \item 在 $\sqrt{2}$ 和 $\sqrt{3}$ 之间，写出一个有理数和一个无理数。
\end{enumerate}


\end{practice}



\begin{answer}


\textbf{平方根与算术平方根 练一练}

\begin{enumerate}
  \item $64$ 的平方根为 $\pm 8$ 。
  \item $\sqrt{121} = 11$ 。
  \item $-\sqrt{36} = -6$ 。
  \item $\sqrt{(-5)^2} = \sqrt{25} = 5$ 。
  \item 错。 $\sqrt{16} = 4$ ，不是 $\pm 4$ 。 $\sqrt{16}$ 是算术平方根，只取非负值。 $16$ 的平方根是 $\pm 4$ 。
  \item 边长为 $\sqrt{200} = 10\sqrt{2}$ 米。
\end{enumerate}


\textbf{立方根 练一练}

\begin{enumerate}
  \item $\sqrt[3]{343} = 7$ 。
  \item $\sqrt[3]{-1000} = -10$ 。
  \item $\sqrt[3]{\frac{27}{8}} = \frac{3}{2}$ 。
  \item $\sqrt[3]{-0.125} = -0.5$ 。
  \item 错。每个实数（包括负数）都有唯一的立方根。
  \item $\sqrt[3]{1728} = 12$ ，棱长为 $12$ 厘米。
\end{enumerate}


\textbf{无理数的发现 练一练}

\begin{enumerate}
  \item $\sqrt{36} = 6$ ，是有理数。
  \item $\sqrt{11}$ 是无理数（ $11$ 不是完全平方数）。
  \item $\sqrt{2}$ （因为 $1 < \sqrt{2} < 2$ ）。
  \item $7^2 = 49 < 50 < 64 = 8^2$ ，所以 $\sqrt{50}$ 的整数部分为 $7$ 。
  \item 错。无限循环小数（如 $0.333\ldots$ ）是有理数，只有无限不循环小数才是无理数。
  \item $\sqrt{8}$ 和 $\pi - 3$ 是无理数。 $0.333\ldots = \frac{1}{3}$ 是有理数， $-\frac{22}{7}$ 是有理数。
\end{enumerate}


\textbf{实数的分类与实数轴 练一练}

\begin{enumerate}
  \item 有理数： $\sqrt{25}$ （ $= 5$ ）、 $\frac{1}{3}$ 、 $0$ 。无理数： $\sqrt{7}$ 、 $-\pi$ 。
  \item $\sqrt{10}$ 与 $3.2$ 比较： $(\sqrt{10})^2 = 10$ ， $3.2^2 = 10.24$ ，因为 $10 < 10.24$ ，所以 $\sqrt{10} < 3.2$ 。
  \item $\sqrt{3} \approx 1.73 < 2$ ，所以 $\sqrt{3} - 2 < 0$ ， $\lvert \sqrt{3} - 2 \rvert = 2 - \sqrt{3}$ 。
  \item $4^2 = 16$ ， $5^2 = 25$ ，所以 $4 < \sqrt{20} < 5$ 。 $4.4^2 = 19.36$ ， $4.5^2 = 20.25$ ，所以 $\sqrt{20} \approx 4.5$ 。
  \item 对。实数与数轴上的点一一对应。
  \item $\sqrt{2} \approx 1.414$ ， $\sqrt{3} \approx 1.732$ 。有理数例如 $1.5$ ；无理数例如 $\sqrt{2.5}$ （ $\approx 1.581$ ）。
\end{enumerate}


\end{answer}
